\DSource{0140}{51}
\hypertarget{AB01-PDF0140-P01}{}
\DLine{AB01-PDF0140-body-L001}{}{}{\raisebox{-0.15pt}{\includegraphics[height=6.1pt]{AB01-PDF0140-G01.png}} his distantiis ei altera tantum inaequalitas additur, quae ab excentrico iuxta Lunae a Sole}
\DLine{AB01-PDF0140-body-L002}{}{}{elongationem oritur.}
\hypertarget{AB01-PDF0140-P02}{}
\DLine{AB01-PDF0140-body-L003}{}{}{\Indent Haec est imago quatuor sphaerarum Lunae, ad demonstrationem geometricam inserviens.}
\par\noindent\begin{minipage}[t]{0.48\FullSourceWidth}\vspace{0pt}\centering\hypertarget{AB01-PDF0140-F01}{}\includegraphics[width=\linewidth]{AB01-PDF0140-F01.png}\end{minipage}\hfill\begin{minipage}[t]{0.48\FullSourceWidth}\vspace{0pt}\setlength{\GutterOffset}{0.52\FullSourceWidth}
\hypertarget{AB01-PDF0140-P03}{}
\DLine{AB01-PDF0140-body-L004}{}{}{\Indent Dicit [auctor]: Circa centrum E sit}
\DLine{AB01-PDF0140-body-L005}{5}{}{circulus pareclipticus (1) ABCD, qui eodem}
\DLine{AB01-PDF0140-body-L006}{}{}{tempore sit etiam circulus deflectens [i. e.}
\DLine{AB01-PDF0140-body-L007}{}{}{obliquus]; similiter sphaerae continget quae}
\DLine{AB01-PDF0140-body-L008}{}{}{circa eius polos revolvitur. Diametrum AS}
\DLine{AB01-PDF0140-body-L009}{}{}{ducamus, in cuius puncto F sit centrum ex-}
\DLine{AB01-PDF0140-body-L010}{10}{}{centrici AMP, radio FA descripti. Arcus AM}
\DLine{AB01-PDF0140-body-L011}{}{p. 78.}{sit motus epicycli a puncto A, apogeo et}
\DLine{AB01-PDF0140-body-L012}{}{}{loco Solis, usque ad M, quod centrum epi-}
\DLine{AB01-PDF0140-body-L013}{}{}{cycli GHRK ponamus. [Punctum F′ in dia-}
\DLine{AB01-PDF0140-body-L014}{}{}{metro ita signetur ut sit EF′ = EF]. Li-}
\DLine{AB01-PDF0140-body-L015}{15}{}{neas EMG, F′MH (2) ducamus; erit G (3)}
\DLine{AB01-PDF0140-body-L016}{}{}{locus apogei in epicyclo, ut centro E Ter-}
\DLine{AB01-PDF0140-body-L017}{}{}{rae et eclipticae apparet; H (4) locus veri}
\DLine{AB01-PDF0140-body-L018}{}{}{apogei qui ex puncto F′ (5) conspicitur. Pa-}
\DLine{AB01-PDF0140-body-L019}{}{}{tet arcum HG esse differentiam [i. e. ae-}
\DLine{AB01-PDF0140-body-L020}{20}{}{quationem] motus anomalistici Lunae in}\end{minipage}\par
\DLine{AB01-PDF0140-body-L021}{}{}{in epicyclo, quae in tertia columna tabularum aequationis Lunae (6) describitur. In epicyclo}
\DLine{AB01-PDF0140-body-L022}{}{}{Luna a G ad H, postea ad R moveatur, denique ad K perveniat. Lineam [E]KN, epicyclum}
\DLine{AB01-PDF0140-body-L023}{}{}{tangentem, ducamus; item MK [perpendicularem], quae erit semidiametrus epicycli [a Sole]}
\DLine{AB01-PDF0140-body-L024}{}{}{declivis quantitate qua centrum epicycli a puncto A excentrici deflectit. Quoniam Luna est}
\DLine{AB01-PDF0140-body-L025}{25}{}{in linea epicyclum tangente, haec semidiametrus epicycli [vel angulus MEK ab eadem sub-}
\DLine{AB01-PDF0140-body-L026}{}{}{tensus] est summa inaequalitatis simplicis et alterius inaequalitatis, quae iuxta elongationem}
\DLine{AB01-PDF0140-body-L027}{}{}{Lunae a puncto A Solis determinatur ({\itshape a}). --- Ut ex hac figura patet, Luna stante in prima}
\DLine{AB01-PDF0140-body-L028}{}{}{medietate GHR epicycli, locus verus Lunae in ecliptica, qui a centro E videtur, minor est}
\DLine{AB01-PDF0140-body-L029}{}{}{loco medio eius in longitudine, qui est centrum epicycli; aequatio igitur de medio motu}
\DLine{AB01-PDF0140-body-L030}{30}{}{Lunae demitur, si anomalia minor quam $180^\circ$ est. Cum autem Luna in altera medietate RKG}
\DLine{AB01-PDF0140-body-L031}{}{}{sit, eius locus verus locum medium in ecliptica superat; qua re, anomalia $180^\circ$ excedente,}
\DLine{AB01-PDF0140-body-L032}{}{}{aequatio, si Deus vult, motui medio Lunae est addenda.}
\hypertarget{AB01-PDF0140-P04}{}
\DLine{AB01-PDF0140-body-L033}{}{}{\Indent Aequatio simplex, quae temporibus syzygiarum apparet et in hoc nostro libro in secunda}
\DLine{AB01-PDF0140-body-L034}{}{}{columna tabularum aequationis (7) describitur, iam diximus quo modo numeris supputetur,}
\DLine{AB01-PDF0140-body-L035}{35}{}{eadem via quam secuti sumus in calculanda et in tabulis describenda aequatione Solis (8).}
\DLine{AB01-PDF0140-body-L036}{}{}{Maxima quantitas, quam haec simplex inaequalitas Lunae attingere potest, est $5^\circ1'$ (9), cuius}
\DLine{AB01-PDF0140-body-L037}{}{}{sinus, qui tunc erit semidiametrus epicycli, $5^p\qfrac{1}{4}$ fere complectitur; talis enim est ratio 60 par-}
\DLine{AB01-PDF0140-body-L038}{}{}{tium semidiametrum efficientium ad $5^p\qfrac{1}{4}$. Id Ptolemaeus (10) declaravit et probavit eclipsibus}
\par\vspace{6pt}\hrule height.25pt\vspace{5pt}
\noindent\begin{minipage}[t]{.48\textwidth}\vspace{0pt}\fontsize{9}{11.5}\selectfont\setlength{\GutterOffset}{0pt}
\hypertarget{AB01-PDF0140-N01}{}
\DLine{AB01-PDF0140-notes_left-L001}{}{}{\Indent (1) Vide supra pag. 50, adn. 1.}
\hypertarget{AB01-PDF0140-N02}{}
\DLine{AB01-PDF0140-notes_left-L002}{}{}{\Indent (2) Cod. et Plato: EMH, FMG.}
\hypertarget{AB01-PDF0140-N03}{}
\DLine{AB01-PDF0140-notes_left-L003}{}{}{\Indent (3) Cod. et Plato H.}
\hypertarget{AB01-PDF0140-N04}{}
\DLine{AB01-PDF0140-notes_left-L004}{}{}{\Indent (4) Cod. et Plato G.}
\hypertarget{AB01-PDF0140-N05}{}
\DLine{AB01-PDF0140-notes_left-L005}{}{}{\Indent (5) Cod. et Plato: « ex centro F scilicet ex cen-}
\DLine{AB01-PDF0140-notes_left-L006}{}{}{tro excentrici ». Punctum F′ apud eos deest.}
\hypertarget{AB01-PDF0140-N06}{}
\DLine{AB01-PDF0140-notes_left-L007}{}{}{\Indent (6) Fol. 189,v. et sqq; sed est tertia columna}
\DLine{AB01-PDF0140-notes_left-L008}{}{}{tantum si columna aequationis Solis non computetur.}
\end{minipage}
\hfill\begin{minipage}[t]{.48\textwidth}\vspace{0pt}\fontsize{9}{11.5}\selectfont\setlength{\GutterOffset}{0pt}
\hypertarget{AB01-PDF0140-N07}{}
\DLine{AB01-PDF0140-notes_right-L001}{}{}{\Indent (7) Fol. 189,v. et seqq.; sed est secunda colu-}
\DLine{AB01-PDF0140-notes_right-L002}{}{}{mna tantum si columna aequationis Solis vel colu-}
\DLine{AB01-PDF0140-notes_right-L003}{}{}{mna numerorum non computetur.}
\hypertarget{AB01-PDF0140-N08}{}
\DLine{AB01-PDF0140-notes_right-L004}{}{}{\Indent (8) Cap. XXVIII.}
\hypertarget{AB01-PDF0140-N09}{}
\DLine{AB01-PDF0140-notes_right-L005}{}{}{\Indent (9) Sunt numeri Ptolemaei.}
\hypertarget{AB01-PDF0140-N10}{}
\DLine{AB01-PDF0140-notes_right-L006}{}{}{\Indent (10) {\itshape Almag.} IV, 5 (ed. Halma, t. I, pag. 244-}
\DLine{AB01-PDF0140-notes_right-L007}{}{}{-260).}
\end{minipage}
